\documentclass[11pt]{article}

\usepackage[margin=1in]{geometry}
\usepackage{amsmath,amssymb}
\usepackage{array}
\usepackage{booktabs}
\usepackage{enumitem}
\usepackage{longtable}
\usepackage{microtype}
\usepackage{xurl}
\usepackage[hidelinks]{hyperref}

\setlist[itemize]{leftmargin=1.5em}
\setlist[enumerate]{leftmargin=1.5em}
\renewcommand{\arraystretch}{1.18}
\newcolumntype{L}[1]{>{\raggedright\arraybackslash}p{#1}}
\emergencystretch=6em
\DeclareUrlCommand{\code}{\urlstyle{tt}}
\newcommand{\filepath}[1]{\path{#1}}

% Keep the main title concise: a 2 to 6 word subject noun phrase.
\newcommand{\dossiertitle}{Short Feature Title}
\newcommand{\dossiersubtitle}{Technical Design Dossier}
\newcommand{\dossierauthor}{Prepared by Codex}
\newcommand{\preparedtimestamp}{Prepared Saturday, 23 May 2026, 00:30:53 EDT (-0400)}
\newcommand{\repo}{Repository or workspace path}
\newcommand{\artifactstatus}{Draft for technical review}
\newcommand{\classification}{Internal technical document}

\hypersetup{pdftitle={\dossiertitle}, pdfauthor={\dossierauthor}}

\title{\dossiertitle\\[0.35em]\large\dossiersubtitle}
\author{\dossierauthor}
\date{\preparedtimestamp}

\begin{document}
\maketitle

\begin{center}
\begin{tabular}{>{\bfseries}p{0.24\linewidth}p{0.62\linewidth}}
Repository & \repo \\
Status & \artifactstatus \\
Classification & \classification \\
Canonical artifact & Replace with relative path to this TeX source \\
Companion PDF & Replace with relative path to rendered PDF \\
\end{tabular}
\end{center}

\begin{abstract}
Replace this paragraph with a concise technical summary: the problem, the selected design, the evidence supporting it, and the validation standard for accepting the work. Keep this section factual and decision-oriented.
\end{abstract}

\tableofcontents
\newpage

\section{Purpose}
State why this dossier exists, who should use it, and what decision or implementation pass it is meant to make unambiguous.

\section{Glossary}
Define specialized terms before relying on them. Omit this section only when the document uses ordinary project vocabulary.

\begin{longtable}{L{0.26\linewidth}L{0.62\linewidth}}
\toprule
\textbf{Term} & \textbf{Definition} \\
\midrule
\endhead
Replace term & Replace with the precise meaning used in this dossier. \\
\bottomrule
\end{longtable}

\section{Current Status and Source Evidence}
Summarize only facts verified from the live repository, current reports, runtime behavior, or user-supplied source material.

\begin{longtable}{L{0.25\linewidth}L{0.35\linewidth}L{0.30\linewidth}}
\toprule
\textbf{Evidence} & \textbf{Observed fact} & \textbf{Design implication} \\
\midrule
\endhead
\filepath{path/to/file.ext} & Replace with verified behavior, interface, constant, or test coverage. & Replace with why this constrains the design. \\
\code{command output} & Replace with relevant result. & Replace with validation or blocker implication. \\
\bottomrule
\end{longtable}

\section{Problem Diagnosis}
Define the failure mode or design gap. Distinguish the visible symptom from the owner, invariant, or contract that actually needs to change.

\section{Goals and Non-Goals}
\subsection{Goals}
\begin{itemize}
  \item Replace with a measurable outcome.
  \item Replace with a repository-visible behavioral or architectural outcome.
\end{itemize}

\subsection{Non-Goals}
\begin{itemize}
  \item Replace with a plausible scope item that is explicitly excluded.
  \item Replace with a compatibility path or fallback that should not be preserved unless a real boundary requires it.
\end{itemize}

\section{Assumptions and Constraints}
\begin{longtable}{L{0.22\linewidth}L{0.42\linewidth}L{0.26\linewidth}}
\toprule
\textbf{Type} & \textbf{Statement} & \textbf{Validation path} \\
\midrule
\endhead
Assumption & Replace with an assumption the implementation depends on. & Replace with command, test, inspection, or owner confirmation. \\
Constraint & Replace with a hard boundary from the repo, runtime, API, or product requirement. & Replace with proof path. \\
\bottomrule
\end{longtable}

\section{Canonical Ownership Model}
Name the invariant, the deepest behavior owner, and the chosen edit layer.

\begin{longtable}{L{0.24\linewidth}L{0.34\linewidth}L{0.32\linewidth}}
\toprule
\textbf{Concern} & \textbf{Canonical owner} & \textbf{Reason} \\
\midrule
\endhead
Replace concern & Replace owner path, module, contract, or process. & Replace with first-principles rationale. \\
\bottomrule
\end{longtable}

\section{Cutover and Compatibility Policy}
Prefer hard cutovers inside the repository's semantic core. Allow compatibility only at verified external boundaries, and record the removal condition.

\begin{longtable}{L{0.24\linewidth}L{0.31\linewidth}L{0.35\linewidth}}
\toprule
\textbf{Surface} & \textbf{Policy} & \textbf{Boundary or removal condition} \\
\midrule
\endhead
Replace core path, field, command, or owner & Hard cutover; remove duplicate semantics in the same change. & Not applicable, or specify final migration proof. \\
Replace external boundary & Temporary translation only for a verified consumer or persisted format. & Replace with version, adoption, telemetry, or deletion proof. \\
\bottomrule
\end{longtable}

\section{Design}
\subsection{Selected Design}
Describe the selected architecture in terms of responsibilities, data flow, state ownership, and boundary contracts. Keep implementation details only where they prevent ambiguity.

\subsection{Contracts}
\begin{longtable}{L{0.22\linewidth}L{0.35\linewidth}L{0.33\linewidth}}
\toprule
\textbf{Contract} & \textbf{Required behavior} & \textbf{Consumer or proof} \\
\midrule
\endhead
Replace contract name & Replace with required behavior, schema, state transition, or numerical invariant. & Replace with consumer, test, or rendered proof. \\
\bottomrule
\end{longtable}

\subsection{Algorithms or Formal Model}
Include equations or pseudocode only where they clarify the design. For example:
\[
  y = f(x, t; \theta)
\]
where \(x\) is the input state, \(t\) is time or generation, and \(\theta\) is the selected parameter set.

\section{Alternatives Considered}
\begin{longtable}{L{0.24\linewidth}L{0.24\linewidth}L{0.42\linewidth}}
\toprule
\textbf{Alternative} & \textbf{Disposition} & \textbf{Reason} \\
\midrule
\endhead
Replace alternative & Accepted or rejected & Replace with the tradeoff that determines the decision. \\
\bottomrule
\end{longtable}

\section{Implementation Plan}
Each phase should include tests before code changes whenever practical.

\begin{enumerate}
  \item \textbf{Phase 1: Failing proof.} Add or identify the test, visual canary, trace, or command that currently fails for the right reason.
  \item \textbf{Phase 2: Canonical owner change.} Modify the deepest owner identified in the Canonical Ownership Model.
  \item \textbf{Phase 3: Boundary cleanup.} Remove stale aliases, duplicate paths, shadow state, or compatibility wrappers unless a verified external boundary requires translation.
  \item \textbf{Phase 4: Validation.} Run the focused proof first, then the broader regression lane appropriate to the repository.
\end{enumerate}

\section{Test-Driven Development Contract}
\begin{longtable}{L{0.26\linewidth}L{0.34\linewidth}L{0.30\linewidth}}
\toprule
\textbf{Test or proof} & \textbf{Expected initial result} & \textbf{Acceptance result} \\
\midrule
\endhead
\code{replace command} & Fails or demonstrates the existing gap. & Passes and proves the selected invariant. \\
\bottomrule
\end{longtable}

\section{Validation and Acceptance Criteria}
\begin{itemize}
  \item Replace with exact command and expected artifact.
  \item Replace with metric, threshold, visual criterion, or behavioral invariant.
  \item Replace with evidence that would falsify the design.
\end{itemize}

\section{Risks and Failure Modes}
\begin{longtable}{L{0.23\linewidth}L{0.34\linewidth}L{0.33\linewidth}}
\toprule
\textbf{Risk} & \textbf{Failure mode} & \textbf{Mitigation or blocker} \\
\midrule
\endhead
Replace risk & Replace with concrete failure mode. & Replace with mitigation, test, or reason to stop. \\
\bottomrule
\end{longtable}

\section{Exact Next Sequence}
\begin{enumerate}
  \item Replace with the first command, file inspection, or test.
  \item Replace with the first code or document edit.
  \item Replace with the focused validation command.
  \item Replace with the broader verification command.
\end{enumerate}

\section{Exit Rule}
The implementation is complete only when the canonical owner has been changed, stale competing paths have been removed or justified at a boundary, the focused proof passes, the broader validation lane passes, and this dossier's companion PDF has been rendered from the current TeX source.

\section{References}
\begin{itemize}
  \item Replace with source files, reports, papers, docs, or command outputs cited in the dossier.
\end{itemize}

\appendix
\section{Appendix: Deferred Evidence}
Place detailed logs, long tables, raw measurements, or source excerpts here when they would interrupt the main argument.

\end{document}
